\section{A worked compilation across the existing languages}
\label{sec:overview}
A compiler overview should expose where a decision is made, what that
choice commits subsequent passes to, and which syntax remains residual.
We therefore follow one executable program through every existing boundary.
The example uses a computed join key and a computed answer component, so its
compilation cannot be explained by a permutation of ordinary CQ atoms alone.
All displayed generated names come from the recorded implementation. Listings
normalize whitespace, and the RA listing omits raw-identifier sigils that denote
the same attribute names; explicitly marked abbreviations are explanatory
notation. Complete generated artifacts and their executable checks are
reproducible with the driver in \code{doc/formal/example}.

\subsection{Source program and observation}
\begin{lstlisting}
pub struct Overview;
relation R(a: i32, b: i32);
relation S(a: i32, b: i32);
relation answer(a: i32, b: i32);
answer(x, *y + 1) :- R(x, k), let next: i32 = *k + 1,
                    S(next, y), if y > 0;
\end{lstlisting}
The declaration of \code{answer} supplies the Rust types of the computed
answer. The first atom binds logical variables \code{x} and \code{k} to
input values. The host expression in \code{let} computes \code{next}; its
later bare occurrence in \code{S(next,y)} is a required equality key, not a
new independent producer of that value. The guard tests \code{y}; the head
computes its second positional component. These occurrences of \code{*k}
and \code{*y} are ordinary Rust dereferences in native expression ASTs, not
new arithmetic forms in the CQ grammar. They reflect the host interface's
borrowed aliases for input-derived fields.

Use the concrete input instance
\[
 I(R)=\{(1,6),(2,6),(3,8)\},\qquad
 I(S)=\{(7,10),(7,-2),(9,4)\}.
\]
For this non-overflowing instance, ordinary integer evaluation gives the
answer relation
\[
 \{(x,y+1)\mid (x,k)\in I(R),\ n=k+1,\ (n,y)\in I(S),\ y>0\}
 =\{(1,11),(2,11),(3,5)\}.
\]
This equation states the intended finite result for the example; it does not
assert that every accepted Rust expression denotes a total mathematical
function. In particular, unrestricted \code{i32} addition would require an
account of overflow before such a universal statement.

There are two public compositions. \code{mini_linq!} selects the default
logical, index, and lazy physical passes. \code{planned_query!} selects the
planned logical pass, exact result-index selection, and materialized physical
pass. On this input both return the sorted vector
\code{[(1,11),(2,11),(3,5)]} through \code{run}. The driver also compiles and
executes each route's fully lowered ordinary Rust file. This is evidence for
this one program and observation, not equivalence for arbitrary host effects.

\subsection{The boundaries and the decisions they own}
\begin{center}
\small
\begin{tabular}{>{\raggedright\arraybackslash}p{.23\linewidth}>{\raggedright\arraybackslash}p{.34\linewidth}>{\raggedright\arraybackslash}p{.35\linewidth}}
\toprule
Existing target & Decision made on this edge & Structure left explicit \\
\midrule
\code{cq::Module} & Parse the DSL shell and native Rust leaves; check source binding premises. & Atoms, expression arguments, host patterns, guards, and head ASTs. \\
\code{relational_plan::}\newline\code{Module} & Choose a legal clause/join schedule; assign named algebra results and their sorts. & Algebra operators and unevaluated host expressions. \\
\code{index_requirements::}\newline\code{Module} & Add required equality keys to the unchanged algebra. & Exact indexed operand names and column positions. \\
\code{syn::File} containing \code{pull!} & Choose storage, row layout, aliases, access expressions, and materialization boundaries. & Existing Access plans quoted inside ordinary Rust. \\
\code{iterator_pipeline::}\newline\code{Pipeline} & Expose the local scan/join/filter/project/distinct sequence and binding order. & Native Rust leaves inside named iterator operators. \\
\code{syn::Expr} & Implement each local operator with concrete Rust adapters and suspension. & Ordinary Rust for the host compiler to resolve and type-check. \\
\bottomrule
\end{tabular}
\end{center}
The exact typed entry points for the two file-level routes are:
\begin{center}
\footnotesize
\begin{tabular}{lll}
\toprule
Edge & Default route & Planned route \\
\midrule
CQ $\to$ RA & \code{compile_cq} & \code{compile_cq_planned} \\
RA $\to$ indexes & \code{compile_relational_plan} & \code{compile_result_indexes} \\
Indexes $\to$ Rust file & \code{compile_index_requirements} & \code{compile_materialized} \\
\bottomrule
\end{tabular}
\end{center}
Both reuse \code{compile_pull} and \code{compile_iterator_pipeline} for
the final two local edges.

The fourth row is not a new file-wide IR: it is the ordinary Rust AST with
existing lower-language macro holes. Its embedded \code{rust_access_plan::Plan}
is the source of \code{compile_pull}. Thus the physical pass returns a
\code{syn::File}, while the last local pass returns a \code{syn::Expr} that
fills one hole in that file. The distinction between the surrounding Rust
context and the embedded language is essential to composition.

\subsection{CQ to named relational algebra: a value dependency}
The dependency judgment records required logical names, provided logical
names, and the current movement classification. For the source above its
facts are:
\begin{center}
\begin{tabular}{cllll}
\toprule
Clause & Form & Required & Provided & Movable \\
\midrule
0 & \code{R(x,k)} & $\emptyset$ & $\{x,k\}$ & yes \\
1 & \code{let next = *k+1} & $\{k\}$ & $\{\mathit{next}\}$ & no \\
2 & \code{S(next,y)} & $\{\mathit{next}\}$ & $\{\mathit{next},y\}$ & yes \\
3 & \code{if y>0} & $\{y\}$ & $\emptyset$ & yes \\
\bottomrule
\end{tabular}
\end{center}
``Provided'' includes a bare atom's equality columns, including an already
bound name; it does not mean every listed name is fresh. In clause 2,
\code{next} is both offered as a relation column and required from its
preceding \code{let}. The computed binding is a movement barrier under the
current conservative classification. Its lack of external state in this
particular expression does not cause the compiler to infer a general
purity theorem.

The recorded \code{Auto} derivation is exactly
\[
 \mathsf{Match}(0);\ \mathsf{Bind}(1);\ \mathsf{Match}(2);\
 \mathsf{Guard}(3);\ \mathsf{Finish}.
\]
There is no GYO analysis or join-tree certificate on this route. The let and
computed head put the complete query outside the present graph-planning
domain, so \code{compile_cq_planned} selects dependency-constrained prefix
rules. \code{compile_cq} selects the written policy for the same extended
source. They happen to produce the same algebra here:
\begin{lstlisting}
relational {
  r0 = rename R { a -> x, b -> k };
  r1 = extend r0 with { next: i32 }
       value({ let next: i32 = *k + 1; (next,) });
  r2 = rename S { a -> next, b -> y };
  r3 = natural_join r1 with r2;
  r4 = select r3 if (y > 0);
  r5 = project r4 keep { k, next, x, y };
  output r5 as answer(x, *y + 1);
}
\end{lstlisting}
The complete target additionally retains the source program and declarations;
they are omitted from this listing. The native Rust expression remains in
\code{Extend}; logical planning does not execute it. Its explicit result
column makes the equality key available to \code{NaturalJoin}. The final
\code{Project} retains all witness columns because the head is computed.
That is the current prefix rule's conservative policy: this example does
not demonstrate unused-column pruning, nor does its guard demonstrate
selection pushdown into the right input. On the restricted graph domain,
separate P-* rules perform eager selection and retention.

A particular valuation satisfying the extended query on this instance is
\[
 \nu(x)=1,\quad\nu(k)=6,\quad\nu(\mathit{next})=7,\quad\nu(y)=10.
\]
After occurrence renaming, its corresponding named tuple in \code{r3} has
sort $\{k,\mathit{next},x,y\}$ and those four attribute values. The named
output tuple of \code{r5} remains that four-attribute object; the head's
ordered answer is $(1,11)$. These are three related but different objects:
a valuation of logical variables, a named tuple over a sort, and a positional
answer tuple. The textbook correspondence is established in the foundations.

The concrete sorts, physical field order, and measured materialized counts
are as follows. All listed field types are \code{i32}; the displayed field
order is the implementation's lexical attribute-name order, not a property
of a mathematical sort.
\begin{center}
\begin{tabular}{clrc}
\toprule
Result & Ordered physical fields & Arity & Number of tuples \\
\midrule
\code{r0} & $(k,x)$ & 2 & 3 \\
\code{r1} & $(k,\mathit{next},x)$ & 3 & 3 \\
\code{r2} & $(\mathit{next},y)$ & 2 & 3 \\
\code{r3} & $(k,\mathit{next},x,y)$ & 4 & 5 \\
\code{r4} & $(k,\mathit{next},x,y)$ & 4 & 3 \\
\code{r5} & $(k,\mathit{next},x,y)$ & 4 & 3 \\
\bottomrule
\end{tabular}
\end{center}
The reported peak arity is 4. This is a column count, not a cardinality
estimate, treewidth, or fractional hypertree width.

\subsection{Algebra to index requirements: name the exact operand}
The intersection of \code{r1}'s and \code{r2}'s sorts is
$\{\mathit{next}\}$. The right operand's ordered fields are
$(\mathit{next},y)$, so its key occupies position 0. Result-index selection
preserves the entire algebra and appends exactly
\begin{lstlisting}
indexes { @r2[0]; }
\end{lstlisting}
The \code{@} identifies a named RA result. The requirement is not a guessed
index on whichever base relation might eventually supply a similar column.
For this extended linear plan the default index entry also takes the
result-index route, so both compositions have the same overlay. No index
is requested for the let alone or for the guard. The index pass neither
chooses a different join order nor decides whether the named results will
be materialized.

The representation chosen by the physical pass is a
\code{HashMap<(i32,), Vec<(i32,)>>}: key fields contain \code{next}; bucket
payload fields contain \code{y}. It represents
\[
 7\mapsto\{(10),(-2)\},\qquad 9\mapsto\{(4)\},
\]
where braces abstract the bucket order inherited from an unspecified
HashSet enumeration. A fixed vector retains that insertion order across
lookups. The tuple syntax in generated Rust is \code{(10,)} and so forth.
All source and renamed rows are deduplicated into \code{HashSet}s before
bucket construction, and every renamed row contributes exactly one payload.
A lookup of key 7 therefore yields two candidates, not one representative.

\subsection{Index requirements to staged Rust: storage and residual Access}
The lazy physical pass fuses this accepted linear algebra into one local
Access plan. \code{OverviewStorage} owns
\begin{lstlisting}
rows0: HashSet<(i32, i32)>, // fields (k,x)
rows2: HashSet<(i32, i32)>, // fields (next,y)
index0: HashMap<(i32,), Vec<(i32,)>>,
\end{lstlisting}
Qualification of standard type paths is abbreviated here. \code{build}
constructs these values. The retained \code{rows2} field is actual current
output, even though \code{query} uses its index for this example. The
physical pass has already chosen Rust identifiers
\[
 v0\leftrightarrow k,\quad v1\leftrightarrow\mathit{next},\quad
 v2\leftrightarrow x,\quad v3\leftrightarrow y.
\]
During lazy traversal, \code{v0}, \code{v2}, and \code{v3} have type
\code{&i32}; computed \code{v1} has type \code{i32}. These choices are
physical elaboration facts, not CQ variable types silently changed by the
logical planner.

The body of the generated \code{query} contains the following Access hole.
This excerpt removes \code{#[allow(unused_variables)]} attributes and
normalizes redundant parentheses; its aliases and expressions are otherwise
those of the generated target.
\begin{lstlisting}
::mini_linq::pull! {
  output(v0,v1,v2,v3) =>
    ((v2.clone(), { let y = v3; *y + 1 },)) :-
  rows(v0,v2) => for (v0,v2,) in (rows0.iter()),
  computed(v1) => for (v1,) in
    (::core::iter::once({
      let k = v0;
      { let next: i32 = *k + 1; (next,) }
    })),
  rows(v1,v3) => for (v3,) in
    (index0.get(&(v1.clone(),)).into_iter().flatten()),
  if (v3 > &0i32) => if (v3 > &0i32);
}
\end{lstlisting}
The left sides are logical occurrence labels. The right sides are concrete
Rust iteration and testing expressions. The dependent let is implemented
by a singleton iterator that contributes one owned computed field per
prefix tuple; the later lookup contributes only its fresh payload field.
The key's value is already available. Access lowering therefore need not
rediscover the join key, choose an index, or rewrite the user's arithmetic.
The \code{pull!} body is quotation of an existing lower language; \code{build},
field ownership, aliases, and the surrounding Rust method are already fixed.

The materialized physical pass makes a different commitment on the same
indexed algebra. It emits \code{relation0} through \code{relation5} as
\code{HashSet}s with the ordered fields in the preceding table, and one
local \code{pull!} hole per named result. A typical boundary has the form
\begin{lstlisting}
let relation1: HashSet<(i32,i32,i32)> = {
  let __left = &relation0;
  Iterator::collect(::mini_linq::pull! { /* r1 Access plan */ })
};
\end{lstlisting}
This is an explicitly abbreviated structural excerpt. The complete artifact
contains the actual hole, which scans \code{__left}, evaluates the singleton
computed source, and emits owned $(k,\mathit{next},x)$ tuples. After building
\code{relation2}, this route constructs the same key/payload index shape;
its join hole scans \code{relation1} and probes that index. Its answer field
is a \code{HashSet<(i32,i32)>}; the final output mapping uses
\code{__row.2} for \code{x} and the alias \code{let y = &__row.3} for the
computed head. All relational work occurs during \code{build}, and its
\code{query} method iterates the stored answer. In contrast, the lazy route
runs the let, lookup, guard, and head as the cursor is polled. Agreement on
this pure example does not erase that operational distinction.

\subsection{Access to Pipeline to ordinary Rust}
For the one-hole lazy route, \code{compile_pull} returns six named operators:
\begin{lstlisting}
iter0 = scan (v0,v2,) in (rows0.iter()) yield (v0,v2,);
iter1 = join iter0 as (v0,v2,) with (v1,) in
  (once({ let k = v0; { let next: i32 = *k+1; (next,) } }))
  yield (v0,v2,v1,);
iter2 = join iter1 as (v0,v2,v1,) with (v3,) in
  (index0.get(&(v1.clone(),)).into_iter().flatten())
  yield (v0,v2,v1,v3,);
iter3 = filter iter2 as (v0,v2,v1,v3,) if (v3 > &0i32);
iter4 = project iter3 as (v0,v2,v1,v3,)
  yield ((v2.clone(), { let y = v3; *y+1 },));
iter5 = distinct iter4;
return iter5;
\end{lstlisting}
The preceding \code{once} abbreviates \code{::core::iter::once}; the same
allow attributes are omitted as above. The payload order is now binding
introduction order $(k,x,\mathit{next},y)$, whereas named RA storage used
$(k,\mathit{next},x,y)$. This change is visible in the operator patterns and
yields. It does not change either the mathematical sort or the meaning of
\code{v1}. Explicit patterns are the correspondence witness for the new
layout.

Each local operator is lowered separately by
\code{compile_iterator_pipeline}: Scan becomes \code{map}; both Join
operators become \code{flat_map} with an inner \code{map}; Filter becomes
consuming \code{filter_map}; Project becomes \code{map}; and Distinct
becomes a \code{filter} with an owned seen \code{HashSet}. In the lookup
operator, the exact generated structure is:
\begin{lstlisting}
let iter2 = ::core::iter::Iterator::flat_map(
  iter1,
  move |(v0,v2,v1)| {
    ::core::iter::Iterator::map(
      ::core::iter::IntoIterator::into_iter(
        index0.get(&(v1.clone(),)).into_iter().flatten()),
      move |(v3,)| {
        let v0 = ::core::clone::Clone::clone(&v0);
        let v2 = ::core::clone::Clone::clone(&v2);
        let v1 = ::core::clone::Clone::clone(&v1);
        (v0,v2,v1,v3)
      })
  });
\end{lstlisting}
Cloning the retained prefix fields occurs inside the inner map, once per
matching payload. This placement is part of the operational lowering, even
when cloning \code{i32} or its references is observationally silent. The
final \code{syn::Expr} wraps the six locals in
\begin{lstlisting}
Iterator::flatten(iter::once_with(move || {
  /* iter0 ... iter5 definitions */
  iter5
}))
\end{lstlisting}
Here standard path qualification and the body are explicitly abbreviated.
The wrapper delays construction until the first poll. No CQ atom or global
join-order decision remains at this edge. The host compiler still performs
name resolution, type checking, borrowing checks, and code generation for
the resulting ordinary Rust.

\subsection{What this example establishes about composition}
The driver invokes \code{compile_cq} or \code{compile_cq_planned}, their
respective index and physical functions, then traverses the returned Rust
AST to expose and lower existing macro holes with \code{compile_pull} and
\code{compile_iterator_pipeline}. It also invokes the two public CQ macros
as actual proc macros. Its twelve generated source files preserve all six
boundaries of both compositions; \code{report.txt} records the derivation,
and \code{validation.txt} records independently asserted result tuples,
intermediate arities/counts, and physical counters. The intermediate counts
are $(3,3,3,5,3,3)$, index rows are 3, and join pairs are 5. These checks make
the walkthrough inspectable without treating a hand-written sketch as
compiler output.

The two logical policies inhabit the same CQ-to-RA edge. The physical
functions also share an indexed-RA source and ordinary staged-Rust target,
but have different supported shapes and operational commitments. A caller
of the typed functions can compose only when the next pass's domain and
required observation relation hold. Shared Rust function types alone do
not establish substitutability. The public macros currently select fixed
continuations; they do not expose a user-configurable CPS parameter.
Invoking \code{pull!} or \code{iterator_pipeline!} directly is an existing
lower-language entry, and quotation embeds that entry in ordinary Rust.
It is not evidence that arbitrary changes of materialization strategy
preserve effect traces.

For comparison, a query consisting of distinct bare-variable atoms, a
normalized head, and admitted primitive guards enters the graph rules.
There, \code{Auto} attempts GYO, performs the two semijoin sweeps if a join
tree is available, and assembles named joins with selection and retention.
This is a distinct algorithm inside the same pass. The present example's
let forces the prefix route for the entire query: graph planning of
separate regions around host-expression barriers is not yet implemented.
Neither route currently supplies negation, aggregation, WCOJ, or cardinality
estimation merely because the parser can represent richer clauses.
